\documentclass[11pt]{article}
\usepackage{amsmath, amssymb}

\title{CSE 400 Lecture Scribe\\
Lecture 21 \& 22: Inequalities and Tail Bounds}
\date{}

\begin{document}
\maketitle

\section{Tail and Motivation of Tail Bounds}

\subsection{Definition of Tail}
The tail of a random variable $X$ is:
\[
P\{X > x\}
\]

\subsection{Why Do We Care About Tails?}
\begin{itemize}
\item Jobs delayed more than 24 hours
\item Packets dropped due to full buffer
\end{itemize}

\textbf{Key Idea:}
\begin{itemize}
\item Mean $\rightarrow$ average behavior (easy)
\item Tail $\rightarrow$ rare extreme events (hard)
\end{itemize}

\subsection{Why Are Tails Hard to Compute?}
\begin{itemize}
\item Requires summing many small probabilities
\end{itemize}

\subsection{Tail Examples}

\textbf{Example 1: Load Balancing}
\begin{itemize}
\item Distribute $n$ jobs among $n$ servers randomly
\item Let $X \sim \text{Binomial}(n,p)$
\item Tail probability:
\[
P\{X > k\} = \sum \binom{n}{i} p^i (1-p)^{n-i}
\]
\end{itemize}

This is a tail event (rare overload).

\textbf{Example 2: Poisson Arrivals}
\begin{itemize}
\item Jobs arrive with rate $\lambda$
\item $X \sim \text{Poisson}(\lambda)$
\item Tail probability:
\[
P\{X > k\} = \sum e^{-\lambda} \frac{\lambda^x}{x!}
\]
\end{itemize}

This represents a tail event (sudden spike in load).

\subsection{Why Tail Bounds?}
\begin{itemize}
\item Exact computation is complicated:
\[
P\{X > k\} = \sum (\text{long sums})
\]
\end{itemize}

\textbf{Idea:}
\begin{itemize}
\item Instead of exact value $\rightarrow$ find an upper bound
\item Easier to compute
\item Still gives a guarantee
\end{itemize}

\subsection{Running Example}
\begin{itemize}
\item $X \sim \text{Binomial}(n, 1/2)$
\item Mean:
\[
E[X] = \frac{n}{2}
\]
\end{itemize}

Question:
\[
P(X > 3n/4)
\]

\section{Markov's Inequality}

\subsection{Definition}
For non-negative random variable $X$:
\[
P(X > a) \leq \frac{E[X]}{a}
\]

\subsection{Key Idea}
\begin{itemize}
\item Uses only the mean
\item Large values are rare
\end{itemize}

\subsection{Proof}
\[
E[X] = \sum x p_X(x)
\]

Consider only $x > a$:
\[
E[X] \geq \sum_{x>a} x p_X(x)
\]

Replace $x$ by $a$:
\[
E[X] \geq \sum_{x>a} a p_X(x) = a P(X > a)
\]

Thus:
\[
P(X > a) \leq \frac{E[X]}{a}
\]

\subsection{Example (Income Distribution)}
\begin{itemize}
\item Average income = 10,000
\item Threshold = 50,000
\end{itemize}

\[
P(X > 50,000) \leq \frac{10,000}{50,000} = 0.2
\]

At most 20\% people can earn 50,000

\subsection{Solved Example 1 (Exam Grades)}
\begin{itemize}
\item $E[X] = 70$, threshold $a = 90$
\end{itemize}

\[
P(X > 90) \leq \frac{70}{90}
\]

Result:
\[
P(X > 90) < 0.778
\]

\subsection{Solved Example 2 (Server Load)}
\begin{itemize}
\item $E[X] = 5000$, $a = 20000$
\end{itemize}

\[
P(X > 20000) \leq \frac{5000}{20000} = 0.25
\]

At most 25\% chance

\section{Chebyshev's Inequality}

\subsection{Definition}
\[
P(|X - \mu| > a) \leq \frac{\text{Var}(X)}{a^2}
\]

\subsection{Key Idea}
\begin{itemize}
\item Uses mean + variance
\item Controls deviation from mean
\end{itemize}

\subsection{Proof}
Apply Markov to $(X - \mu)^2$:
\[
P(|X - \mu| > a) = P((X - \mu)^2 > a^2)
\]
\[
\leq \frac{E[(X - \mu)^2]}{a^2} = \frac{\text{Var}(X)}{a^2}
\]

\subsection{Running Example}
\begin{itemize}
\item $X \sim \text{Binomial}(n, 1/2)$
\end{itemize}

\[
P(X > 3n/4) \leq P(|X - n/2| > n/4)
\]

Bound decreases as $n$ increases

\subsection{Real-Life Example}
\begin{itemize}
\item Mean = 70, std = 10
\item Deviation = 30
\end{itemize}

\[
P(|X - 70| > 30) \leq \frac{10^2}{30^2} = \frac{100}{900} = \frac{1}{9}
\]

At most $\frac{1}{9}$ students are far from mean

\subsection{Solved Example}
\begin{itemize}
\item $\mu = 70$, $\sigma = 5 \Rightarrow \sigma^2 = 25$
\item $X > 90 \Rightarrow |X - 70| > 20$
\end{itemize}

\[
P(X > 90) \leq \frac{25}{20^2}
\]

\section{Chernoff Bound and Poisson Tail}

\subsection{Key Idea}
\begin{itemize}
\item Uses exponential transformation
\item Stronger than Markov and Chebyshev
\item Exponential decay of tail
\end{itemize}

\subsection{Core Trick}
\[
P(X > a) = P(e^{tX} > e^{ta})
\]

Apply Markov:
\[
P(X > a) \leq \frac{E[e^{tX}]}{e^{ta}}
\]

\subsection{Final Form}
\[
P(X > a) \leq \min_{t>0} \frac{E[e^{tX}]}{e^{ta}}
\]

\subsection{Lower Tail}
\[
P(X < a) \leq \min_{t<0} \frac{E[e^{tX}]}{e^{ta}}
\]

\subsection{Poisson Tail Bound}
For $X \sim \text{Poisson}(\lambda)$:
\[
E[e^{tX}] = e^{\lambda(e^t - 1)}
\]

\[
P(X > a) < e^{\lambda(e^t - 1) - at}
\]

\subsection{Solved Example 1 (Binomial Lower Tail)}
\begin{itemize}
\item $X \sim \text{Binomial}(20, 0.5)$, $\mu = 10$
\item Want $P(X < 0)$
\end{itemize}

Steps:
\begin{enumerate}
\item $0 = (1 - \delta)\mu \Rightarrow \delta = 1$
\item Use bound:
\[
P(X < (1 - \delta)\mu) < e^{-\mu \delta^2 / 2}
\]
\item Substitute:
\[
P(X < 0) < e^{-10/2} = e^{-5}
\]
\end{enumerate}

Result:
\[
P(X < 0) < 0.0067
\]

\subsection{Solved Example 2 (Derivation)}
Steps:
\begin{enumerate}
\item Choose $t < 0$:
\[
P(X < a) = P(tX > ta)
\]
\item Exponentiate:
\[
P(e^{tX} > e^{ta})
\]
\item Apply Markov:
\[
\leq \frac{E[e^{tX}]}{e^{ta}}
\]
\item Minimize over $t < 0$
\end{enumerate}

\section{Jensen's Inequality}

\subsection{Convex Functions}
A function $g$ is convex if:
\[
g''(x) > 0
\]

\subsection{Statement}
For convex $g$:
\[
g(E[X]) \leq E[g(X)]
\]

\subsection{Example (Verification)}
Given:
\begin{itemize}
\item $X = 2$ w.p. 0.5, $X = 4$ w.p. 0.5
\item $g(x) = x^2$
\end{itemize}

Steps:
\begin{enumerate}
\item $E[X] = 3$
\item $g(E[X]) = 9$
\item $E[g(X)] = 0.5(4) + 0.5(16) = 10$
\end{enumerate}

Result:
\[
g(E[X]) = 9 < 10 = E[g(X)]
\]

\subsection{Application}
\[
E[X^a] \geq (E[X])^a
\]
for $a > 1$

\section{Case Study and System-Level Understanding}

\subsection{Problem Setting}
\begin{itemize}
\item Demand $X$: requests/sec
\item Capacity $C$
\item Failure event:
\[
X > C
\]
\end{itemize}

Requirement:
\[
P(X > C) < 0.01
\]

\subsection{Applying Inequalities}

\textbf{(a) Markov}
\begin{itemize}
\item Only mean known ($\mu = 10,000$)
\end{itemize}

Result:
\[
C = 1,000,000
\]

\textbf{(b) Chebyshev}
\begin{itemize}
\item Mean = 10,000, std = 2,000
\end{itemize}

\[
\frac{2000^2}{(C - 10000)^2} < 0.01
\]

Result:
\[
C = 30,000
\]

\textbf{(c) Chernoff}
\begin{itemize}
\item Full distribution known (Poisson)
\end{itemize}

\[
P(X > C) < e^{\cdots}
\]

Solve:
\[
C = 14,500
\]

\subsection{Final Insight}

\begin{tabular}{|c|c|c|}
\hline
Inequality & Information Used & Capacity \\
\hline
Markov & Mean only & 1,000,000 \\
Chebyshev & Mean + Variance & 30,000 \\
Chernoff & Full distribution & 14,500 \\
\hline
\end{tabular}

\subsection{Golden Rule}
More statistical information = Tighter bounds = Lower hardware costs

\end{document}